% UsbDevice class documentation — PUP EN24C T2R4 Linux port
%
% Self-contained: compile with pdflatex usbdevice_doc.tex
% To include in a larger report: copy everything between \begin{document}
% and \end{document} (and add the \usepackage lines to your preamble).
%
\documentclass[a4paper, 11pt]{article}

\usepackage[margin=2.5cm]{geometry}
\usepackage{tikz}
\usetikzlibrary{arrows.meta, positioning, fit, backgrounds}
\usepackage{booktabs}
\usepackage{tabularx}
\usepackage{xcolor}
\usepackage{array}
\usepackage{caption}
\usepackage{parskip}
\usepackage[T1]{fontenc}
\usepackage{lmodern}

% --- colour palette ---
\definecolor{clshead}{RGB}{44,62,80}    % dark blue-grey for concrete class header
\definecolor{clsbody}{RGB}{236,240,241} % light grey for class body
\definecolor{ifacehead}{RGB}{41,128,185}% blue for interface header

\begin{document}

% ============================================================
\section*{USB Interface Layer --- \texttt{UsbDevice} Class}
% ============================================================

\subsection*{Overview}

The \texttt{UsbDevice} class is the software component that directly manages
the physical USB connection between the host computer (Raspberry Pi 4/5,
running Linux) and the \textbf{PUP EN24C T2R4 radar module}.

Physically, a \textbf{USB 3.0 cable} connects the two units.
The radar contains a \textbf{Cypress USB microcontroller} (acting as the
USB device controller); on the host side, \texttt{UsbDevice} drives it using
the open-source \textbf{libusb-1.0} library---no proprietary Windows driver
is required.

Two USB transfer modes are used:
\begin{itemize}
  \item \textbf{Control transfers} --- small, low-rate command/response packets
        used to send radar configuration (carrier frequency, sweep bandwidth,
        transmitter/receiver selection, etc.).
  \item \textbf{Bulk transfers} --- high-throughput, unidirectional streams used
        to push IQ (in-phase/quadrature) sample data continuously from the
        radar to the host during a capture.
\end{itemize}

\bigskip

% ============================================================
\subsection*{Class Structure (UML)}
% ============================================================

\begin{figure}[h!]
\centering
\begin{tikzpicture}[
    font=\small\ttfamily,
    % Style for a header cell (coloured bar at top of each UML compartment)
    hdr/.style={
        draw, rectangle,
        minimum width=7.2cm, text width=7.0cm,
        inner sep=5pt, align=center,
        fill=clshead, text=white,
        font=\small\bfseries\ttfamily
    },
    % Blue header for the interface
    ihdr/.style={hdr, fill=ifacehead},
    % Style for a body compartment (methods / fields list)
    body/.style={
        draw, rectangle,
        minimum width=7.2cm, text width=6.8cm,
        inner sep=5pt, align=left,
        fill=clsbody
    },
]

% ---- IUsbBackend interface ----
\node[ihdr] (ih)
    {\textit{\guillemotleft interface\guillemotright} \quad IUsbBackend};
\node[body, below=0pt of ih] (ib) {
    + open(vid : uint16, pid : uint16) \\[1pt]
    + close() \\[1pt]
    + isOpen() : bool \\[1pt]
    + claimInterface(iface : int) \\[1pt]
    + setAltSetting(iface : int, alt : int) \\[1pt]
    + controlTransfer(bmRequestType, bRequest, \\
    \quad\quad wValue, wIndex, data, len, timeout) \\[1pt]
    + bulkWrite(endpoint, data, len, timeout) \\[1pt]
    + bulkRead(endpoint, buf, len, timeout) \\[1pt]
    + reopenAfterRenumeration(candidates, n, timeout)
};

% ---- UsbDevice concrete class ----
\node[hdr, below=2.2cm of ib] (uh) {UsbDevice};
\node[body, below=0pt of uh] (uf) {
    \textit{--- private fields ---} \\[2pt]
    -- ctx\_ : libusb\_context* \\[1pt]
    -- handle\_ : libusb\_device\_handle* \\[1pt]
    -- claimed\_iface\_ : int
};
\node[body, below=0pt of uf] (um) {
    \textit{--- public methods ---} \\[2pt]
    + UsbDevice() \quad\textit{// initialises libusb context} \\[1pt]
    + \textasciitilde UsbDevice() \quad\textit{// auto-closes on destruction} \\[2pt]
    \textit{(implements all IUsbBackend methods)}
};

% Dashed open-triangle arrow = "implements interface"
\draw[
    -{Triangle[open, length=9pt, width=9pt]},
    dashed, thick
] (uh.north) -- (ib.south)
    node[midway, right=4pt, font=\footnotesize\rmfamily\itshape, text=gray]
    {implements};

\end{tikzpicture}
\caption{UML class diagram. \texttt{UsbDevice} is the concrete
         \texttt{libusb-1.0} implementation of the \texttt{IUsbBackend}
         interface. A test-double (\texttt{FakeUsbBackend}) can substitute
         at the same interface boundary without touching hardware.}
\end{figure}

\clearpage

% ============================================================
\subsection*{Method Reference}
% ============================================================

\begin{table}[h!]
\centering
\small
\renewcommand{\arraystretch}{1.45}
\begin{tabularx}{\textwidth}{
    >{\ttfamily\raggedright}p{4.4cm}
    >{\raggedright\arraybackslash}X
    >{\raggedright\arraybackslash}p{2.6cm}
}
\toprule
\normalfont\textbf{Method} &
\textbf{What it does} &
\textbf{USB mechanism} \\
\midrule

open(vid, pid)
& Scans all USB devices visible to the OS and opens the first one whose
  Vendor ID and Product ID match. Before firmware is loaded, the radar
  appears as the Cypress boot-loader (VID\,\texttt{0x04B4},
  PID\,\texttt{0x8613}).
& Device enumeration \\

close()
& Releases the claimed interface, closes the device handle. Safe to call
  on an already-closed device (no-op).
& --- \\

isOpen()
& Returns \texttt{true} if a device handle is currently held.
& --- \\

claimInterface(n)
& Takes exclusive ownership of USB interface \textit{n}.  On Linux,
  automatically detaches any kernel driver that is already holding the
  interface. Must be called before any data transfer.
& Interface claim \\

setAltSetting(iface, alt)
& Selects alternate setting \textit{alt} on interface \textit{iface}.
  The Cypress chip exposes the high-throughput bulk data endpoints only
  on alternate setting~1; alternate setting~0 carries no endpoints
  (zero bandwidth).
& Alternate setting \\

controlTransfer(\ldots)
& Issues a USB control transfer: sends a short vendor-defined command to
  the radar or reads a short response. Used for all radar configuration
  packets (PLL registers, modulation, sweep timing, TX/RX enable masks)
  and for the firmware-load request (request code \texttt{0xA0}).
& USB Control transfer \\

bulkWrite(ep, data, len)
& Sends a block of bytes to the radar on bulk OUT endpoint \textit{ep}
  (hardware endpoint \texttt{EP2}, address \texttt{0x02}).
  Radar configuration commands are packed into 512-byte packets and
  delivered here.
& USB Bulk OUT \\

bulkRead(ep, buf, len)
& Reads a block of IQ sample data from the radar on bulk IN endpoint
  \textit{ep} (hardware endpoint \texttt{EP6}, address \texttt{0x86}).
  This is the main high-rate data path during capture. A partial read
  on timeout returns whatever bytes arrived rather than raising an error,
  so the capture loop can retry gracefully.
& USB Bulk IN \\

reopenAfter\-Renumeration
& After firmware is downloaded, the Cypress chip resets and re-appears
  on the bus with a new USB address (and possibly a different PID).
  This method releases the old handle and polls every 100\,ms until the
  device is visible again, then re-opens it. Raises an error if the
  timeout elapses.
& Re-enumeration \\

\bottomrule
\end{tabularx}
\caption{Public API of \texttt{UsbDevice}. Methods are called in the order
         shown in the lifecycle diagram (Figure~\ref{fig:lifecycle}).}
\end{table}

\clearpage

% ============================================================
\subsection*{Connection Lifecycle}
% ============================================================

\begin{figure}[h!]
\centering
\begin{tikzpicture}[
    node distance=0.6cm,
    font=\small,
    % Main flow steps
    step/.style={
        draw, rectangle, rounded corners=4pt,
        minimum width=8.5cm, minimum height=0.82cm,
        align=center, fill=clsbody, thick
    },
    % Side annotations
    ann/.style={
        font=\footnotesize\itshape, text=gray,
        text width=4.8cm, align=left
    },
    % Arrow style
    arr/.style={-{Latex[length=7pt, width=5pt]}, thick},
    % Horizontal leader to annotation
    ldr/.style={densely dotted, gray}
]

% --- nodes ---
\node[step] (plug)
    {Radar plugged in via USB 3.0};

\node[step, below=of plug] (open)
    {\texttt{open(0x04B4, 0x8613)} \quad find Cypress boot-loader device};

\node[step, below=of open] (fw)
    {Download firmware into chip RAM via \texttt{controlTransfer} (0xA0 protocol)};

\node[step, below=of fw] (reopen)
    {\texttt{reopenAfterRenumeration} \quad wait for chip to restart};

\node[step, below=of reopen] (claim)
    {\texttt{claimInterface(0)} \quad take exclusive USB access};

\node[step, below=of claim] (alt)
    {\texttt{setAltSetting(0, 1)} \quad activate bulk data endpoints};

\node[step, below=of alt] (cfg)
    {\texttt{controlTransfer} + \texttt{bulkWrite} \quad configure PLL / waveform / TX--RX mask};

\node[step, below=of cfg] (loop)
    {\texttt{bulkRead} loop on EP6 IN \quad stream IQ samples to host};

\node[step, below=of loop] (close)
    {\texttt{close()} \quad release interface and device handle};

% --- arrows ---
\foreach \a/\b in {plug/open, open/fw, fw/reopen, reopen/claim,
                   claim/alt, alt/cfg, cfg/loop, loop/close}
    \draw[arr] (\a) -- (\b);

% Self-loop on bulkRead
\draw[arr, rounded corners=6pt]
    (loop.east) -- ++(0.6,0) -- ++(0,0.62) -- ++(-0.6,0);
\node[font=\scriptsize\itshape, text=gray, right=0.65cm of loop]
    {repeat};

% --- annotations ---
\draw[ldr] (fw.east) -- ++(0.3,0)
    node[ann, right] at ($(fw.east)+(0.3,0)$)
    {The chip has no flash storage; firmware is loaded into RAM at every power-up.};

\draw[ldr] (reopen.east) -- ++(0.3,0)
    node[ann, right] at ($(reopen.east)+(0.3,0)$)
    {Chip disconnects and reconnects with new USB address (re-enumeration).};

\draw[ldr] (loop.east) -- ++(1.4,0)
    node[ann, right] at ($(loop.east)+(1.4,0)$)
    {Continues until capture duration elapses; data saved as interleaved int16 \texttt{.bin}.};

\end{tikzpicture}
\caption{Sequence of \texttt{UsbDevice} operations from hardware plug-in to IQ
         data capture. The firmware-download and re-enumeration steps are the
         most hardware-specific; all subsequent steps are standard USB
         host-side operations.}
\label{fig:lifecycle}
\end{figure}

\end{document}
